% --------------------------------------------------------------------------
\section{Uniform-Anchored Training Objective}
\label{sec:ch4_anchor_objective}

A pure Cluster-DRO objective can become unstable when the cluster weights
\(\mathbf{q}\) concentrate too strongly on a small number of hard clusters. We refer
to this failure mode as \emph{\(q\)-collapse}. In the extreme case, the objective
approaches a worst-cluster loss, where most of the training signal is dominated by
one cluster. This behavior is undesirable in multi-attack adversarial training because
the hardest cluster may change after each cluster refresh, and over-emphasizing one
cluster can reduce robustness on the remaining attacks or groups.

To stabilize training, AttackDRO++ uses a uniform-anchored Cluster-DRO objective
that interpolates between the uniform multi-attack ERM loss and the adaptive
cluster-DRO loss:
\begin{equation}
    \mathcal{L}_{\mathrm{AttackDRO++}}
        =
        (1 - \lambda_{\mathrm{DRO}})\mathcal{L}_{\mathrm{ERM}}
        +
        \lambda_{\mathrm{DRO}}\mathcal{L}_{\mathrm{ClusterDRO}} .
    \label{eq:final_loss}
\end{equation}
Here, \(\lambda_{\mathrm{DRO}}\in[0,1]\) controls the strength of the Cluster-DRO
correction. When \(\lambda_{\mathrm{DRO}}=0\), the method reduces to uniform
multi-attack ERM. When \(\lambda_{\mathrm{DRO}}=1\), the method becomes pure
Cluster-DRO.

Equivalently, Equation~\eqref{eq:final_loss} can be written as
\begin{equation}
    \mathcal{L}_{\mathrm{AttackDRO++}}
        =
        \mathcal{L}_{\mathrm{ERM}}
        +
        \lambda_{\mathrm{DRO}}
        \left(
            \mathcal{L}_{\mathrm{ClusterDRO}}
            -
            \mathcal{L}_{\mathrm{ERM}}
        \right).
\end{equation}
This form shows that \(\lambda_{\mathrm{DRO}}\) scales the difference between the
cluster-reweighted loss and the uniform loss. If the cluster-DRO loss is close to the
uniform loss, the correction is small. If some clusters are substantially harder, the
correction becomes larger, but its influence remains controlled by
\(\lambda_{\mathrm{DRO}}\).

In the main experiments, we use \(\lambda_{\mathrm{DRO}}=0.35\), used as the default anchor strength in the main configuration. This value keeps the uniform multi-attack objective as
the dominant stabilizing term while still allowing the model to adapt toward harder
clusters. Intuitively, this suggests an inverted-U behavior: a small amount of
Cluster-DRO pressure can improve robustness by emphasizing hard groups, whereas
excessive pressure may overfit unstable clusters and harm the balance across attacks.

A second stabilization mechanism is the floor projection applied to the cluster
weights. After each exponentiated-gradient update, \(\mathbf{q}\) is normalized and
projected so that every cluster keeps at least \(q_{\min}\) mass:
\begin{equation}
    q_k \ge q_{\min}, \qquad \sum_{k=1}^{K} q_k = 1.
    \label{eq:floor}
\end{equation}

The anchor controls how much the adaptive cluster loss can influence the total objective,
while the floor projection controls the cluster weights themselves. Together, they
prevent the adaptive component from collapsing onto a small number of clusters and keep
all discovered groups represented during training. In our default setting,
\(q_{\min}=0.03\).

Finally, the \(q\)-update is delayed during the early training phase. For the first
\(T_w\) epochs, \(\mathbf{q}\) remains uniform while the model representation and
initial cluster assignments become more reliable. After this warmup period, the
cluster weights are updated using the exponentiated-gradient rule. In the main
configuration, \(T_w=3\) epochs.
